\section{Solving Linear Equations and Interpolation}

\subsection{Vandermonde system}
We solve the linear system \(Ax=b\) where \(A\in\mathbb{R}^{n\times n}\) is a Vandermonde-type matrix with entries
\[
  A_{ij} = i^{\,n-j}, \qquad i,j\in\{1,\dots,n\},
\]
and \(b\in\mathbb{R}^{n}\) has entries
\[
  b_i = \sum_{k=0}^{n-1} i^k
  = \begin{cases}
  n, & i=1,\\
  \dfrac{i^n - 1}{i-1}, & i>1.
  \end{cases}
\]
The true solution is \(x=\mathbf{1}\). We report the conditioning of \(A\) and the numerical errors \(\lVert x-\mathbf{1}\rVert_2\) and \(\lVert Ax-b\rVert_2\) for \(n\in\{5,10,15\}\).

\begin{table}[H]
\centering
\caption{Vandermonde solve diagnostics}
\label{tab:pset1-vandermonde}
\begin{tabular}{rccc}
\toprule
\(n\) & \(\kappa_2(A)\) & \(\lVert x-\mathbf{1}\rVert_2\) & \(\lVert Ax-b\rVert_2\) \\
\midrule
5 & \(2.62\times 10^{4}\) & \(3.20\times 10^{-13}\) & \(1.14\times 10^{-13}\) \\
10 & \(2.11\times 10^{12}\) & \(7.89\times 10^{-7}\) & \(1.81\times 10^{-7}\) \\
15 & \(2.58\times 10^{21}\) & \(2.17\times 10^{2}\) & \(2.29\) \\
\bottomrule
\end{tabular}
\end{table}

\noindent where \(\kappa_2(A)\) is the condition number of \(A\) in the 2-norm, which is defined as:
\begin{align}
\kappa_2(A) = \|A\|_2 \|A^{-1}\|_2,
\end{align}
where $ \|A\|_2$ is the spectral norm of $A$, corresponding to the largest singular value of $A$. From Table~\ref{tab:pset1-vandermonde}, we see that the condition number grows rapidly as $n$ increases, indicating that the Vandermonde system is ill-conditioned. Hence, for $n$ values larger than 10, the solution is not accurate as columns become nearly linearly dependent in finite precision arithmetic.


\subsection{Polynomial interpolation of \(\log(s)\)}
We approximate \(f(s)=\log(s)\) on \(s\in[1,10]\) by a polynomial of degree \(n-1\),
\[
  \hat f_n(s)=x_1 + x_2 s + \dots + x_n s^{n-1}.
\]
We fit coefficients by (overdetermined) least squares using \(m=100\) equally-spaced points. For \(n\in\{5,10,15,20\}\) we plot the approximant on a fine grid and the approximation error \(\hat f_n(s)-\log(s)\), and we report the Euclidean norm of the error on the fine grid and the condition number of the regression matrix.

\begin{figure}[H]
\centering
\begin{tikzpicture}
\begin{groupplot}[
  group style={group size=2 by 1, horizontal sep=1.4cm},
  width=0.48\textwidth,
  height=0.34\textwidth,
  xlabel={$s$},
  legend style={font=\footnotesize, cells={anchor=west}},
]
\nextgroupplot[
  title={Polynomial approximations},
  ylabel={$f(s)$},
]
\addplot[black, thick] table[x=s,y=true,col sep=comma]{problem1_poly_approx.csv};
\addlegendentry{true}
\addplot[blue] table[x=s,y=pred5,col sep=comma]{problem1_poly_approx.csv};
\addlegendentry{$n=5$}
\addplot[red] table[x=s,y=pred10,col sep=comma]{problem1_poly_approx.csv};
\addlegendentry{$n=10$}
\addplot[green!50!black] table[x=s,y=pred15,col sep=comma]{problem1_poly_approx.csv};
\addlegendentry{$n=15$}
\addplot[purple] table[x=s,y=pred20,col sep=comma]{problem1_poly_approx.csv};
\addlegendentry{$n=20$}

\nextgroupplot[
  title={Approximation error},
  ylabel={$\hat f_n(s)-\log(s)$},
]
\addplot[black, thin] coordinates {(1,0) (10,0)};
\addplot[blue] table[x=s,y=err5,col sep=comma]{problem1_poly_approx.csv};
\addplot[red] table[x=s,y=err10,col sep=comma]{problem1_poly_approx.csv};
\addplot[green!50!black] table[x=s,y=err15,col sep=comma]{problem1_poly_approx.csv};
\addplot[purple] table[x=s,y=err20,col sep=comma]{problem1_poly_approx.csv};
\end{groupplot}
\end{tikzpicture}
\caption{Polynomial approximations of $\log(s)$ (left) and approximation errors (right) for degrees $n-1$ with $n\in\{5,10,15,20\}$, using a monomial basis and least squares on $m=100$ grid points.}
\label{fig:pset1-poly-approx}
\end{figure}
\vspace{1.5em}

\noindent The key numerical lesson from Problem 1.1 carries over directly to polynomial approximation: the least-squares design matrix has columns $1,s,s^2,\dots,s^{n-1}$, so it is \emph{Vandermonde-like} in the monomial basis. As $n$ grows, the columns become increasingly ill-conditioned (similar to Table~\ref{tab:pset1-vandermonde}), so the fitted coefficients can become extremely sensitive to rounding error and to small perturbations in the data. Therefore, increasing $n$ does not guarantee a uniformly better approximation: the function space is richer, but numerical instability can dominate and produce oscillatory or distorted fits (often most visible near the boundaries of $[1,10]$) and non-monotone error behavior.



